\scanpage{69}

donc
\[
\operatorname{Hom}(A_*^n,C_{*M}^{n!})
=\{A_*^n\xrightarrow{\unsure{f^n|f^*}}C_{*M}^n\mid
\sigma^*f=0\ \forall\sigma:\Delta_n\longrightarrow\Delta_m\
\unsure{\text{surj.}}\ m<n\}
\]
\margin{\diagram{\begin{tikzcd}
A_*^n \arrow[r,dash,"f^n"] \arrow[d,"\sigma_A^*"'] & C_*^n) \arrow[d,"\sigma_C^*"]\\
A_*^m \arrow[r] & C_*^m
\end{tikzcd}}}
Or $A_*^n\text{---}C_*^n$ est donné par \deleted{\unsure{$A$}}
\[
A_n^n\xrightarrow{\unsure{K_n|k_n}}M
\]
son composé $A_*^n\text{---}C_*^n\text{---}C_*^m$ est donné par le
\[
A_m^n\longrightarrow M
\]
qui n'est autre que le composé
\[
A_m^n\xrightarrow{\sigma^*}A_n^n\xrightarrow{\unsure{K_n|k_n}}M
\]
et la condition \deleted{\unsure{$f(A_*^n)$}} on a donc
\[
f^n(A_*^n)\subset C_*^n
\Longleftrightarrow
\left|\begin{array}{l}
\unsure{K_n|k_n}\circ\sigma^*=0\\
\forall\sigma:\Delta_n\longrightarrow\Delta_m,\quad m<n
\end{array}\right.
\]
(il suffit de le vérifier qd $\sigma$ est un op de dégénérescence); or cette condition est tjs satisfaite si $A_m^n=0$ si $m<n$, ce qu'on a supposé dans a).

\uline{Ex} Soit $S$ un anneau à \unsure{idéal} $J$ à p.d., $t\in J$, considérons alors cet anneau différentiel $1/2$ simplicial à p.d.
\[
A_*^*=(\Gamma^*(\unsure{\Phi|Q}_{*S})\otimes_S\Lambda^*\Psi_{*S})/(T-t)_{\mathrm{pd}},
\]
c'est une \uline{résolution} de $S$, et $A_j^i=0$ si $i>j$, donc on a un syst. d'\unsure{intégrales|intégrale} de Stokes-DeRham à valeurs dans $S$, donc
\diagram{\begin{tikzcd}
A_*^* \arrow[r,"\sim"]&
\left(\Gamma^*(\unsure{\Phi|Q}_{*S})\otimes_S
\Lambda^*\Psi_{*S}/(T-t)_{\mathrm{pd}}\right)
\arrow[dr] &\\
&&\Lambda^{*+1}\unsure{\Phi|Q}_{*S}
\end{tikzcd}}
dont il faudrait étudier les propriétés \unsure{multiplicatives}.
\margin{\deleted{\unsure{p.d.}}}
\margin{68}
